\subsection{Experiment Setups}
\subsubsection{Platforms and toolchain}

SoCPilot uses DeepSeek~\cite{DeepSeek} as the default LLM, Vivado HLS for accelerator synthesis, Sniper~\cite{sniper} for CPU performance estimation, HyperMapper~\cite{9124618} for design space exploration, and ESP~\cite{mantovani2020agile} for SoC integration and FPGA prototyping. Table~\ref{tab:experimental_setup} summarizes the common experimental environment and reproducibility settings used throughout the evaluation.

%这个地方嵌入式代码手动实现需要强调吗
```latex
\begin{table}[t]
\centering
\caption{Experimental platforms and common evaluation settings.}
\label{tab:experimental_setup}
\renewcommand{\arraystretch}{1.10}
\setlength{\tabcolsep}{4pt}

\begin{tabularx}{\columnwidth}{p{0.30\columnwidth}X}
\toprule
\textbf{Item} & \textbf{Configuration} \\
\midrule

\multicolumn{2}{l}{\textit{LLM and agent configuration}} \\
LLM &
DeepSeek API with \texttt{deepseek-v4-pro}; temperature $=0$. \\

\midrule
\multicolumn{2}{l}{\textit{Processor and software environment}} \\
CPU candidates &
Ariane/CVA6 (RV64), Ibex (RV32), and Leon3 (SPARC V8). \\

Compilation &
RISC-V GCC for Ariane/CVA6 and Ibex; [SPARC toolchain] for Leon3.
CPU-only and CPU--accelerator implementations use identical compiler
settings for each application. \\

Embedded runtime &
Benchmark-specific accelerator-control and runtime code is manually
implemented using the ESP programming model and is held fixed during
evaluation. \\

\midrule
\multicolumn{2}{l}{\textit{Accelerator generation and estimation}} \\
HLS &
Vitis/Vivado HLS 2021.2 with a 10-ns target clock. \\

CPU estimation &
Sniper-based CPU performance estimation; Spike is used for RISC-V
functional simulation where applicable. \\

Area estimation &
A common analytical SoC-area model is used for hardware/software
partitioning and DSE under identical technology assumptions. \\

\midrule
\multicolumn{2}{l}{\textit{System-level exploration and validation}} \\
SoC generation &
ESP-based CPU--accelerator SoC configuration and integration. \\

DSE &
HyperMapper-based exploration over processor, cache, and accelerator
configuration parameters under application-specific area constraints. \\

FPGA validation &
Xilinx Zynq UltraScale+ \texttt{xczu9eg-ffvb1156-2-e};
ESP runtime clock of 75~MHz. \\

\bottomrule
\end{tabularx}
\end{table}
```


% TODO: The final version should explicitly document how CPU, cache/system, and accelerator
% components contribute to the reported SoC-area estimate. If the current model covers only
% accelerator area, RQ3 should instead be phrased as an accelerator-area constraint.

\subsubsection{Baselines and measurement protocol}

Our evaluation follows the progression of the three research questions: RQ1 determines whether the generated system works end to end and improves latency, RQ2 tests whether the specialized workflow is more reliable than open-ended general-agent execution on the hardware/SoC stages common to both methods, and RQ3 evaluates architecture selection under different area constraints.

\textbf{Paired CPU-only baseline.}
For RQ1, each generated SoC is compared against an application-specific paired CPU-only baseline. After SoCPilot selects the CPU type and memory-system parameters, the matched CPU-only configuration retains the same CPU, cache configuration, clock target, compiler, input data, and application algorithm, but executes the selected hot functions on the CPU rather than invoking accelerators. The benchmark-specific embedded code for the accelerated system is manually prepared using the ESP runtime API and is not part of the LLM-generated output. This paired comparison isolates the end-to-end latency benefit of accelerator integration from differences in CPU choice, cache configuration, compiler settings, or application inputs.

Reported latency covers the complete application, including CPU execution, accelerator invocation, data transfer, synchronization, and all unaccelerated functions. Correctness is checked against the CPU reference output before performance results are accepted.

\textbf{General-purpose agent baseline.}
For RQ2, we compare SoCPilot with a general-purpose LLM agent following the HSCO-Bench \cite{hscobench2026} evaluation protocol. Because SoCPilot does not generate benchmark-specific embedded runtime software, the comparison focuses on the automated hardware/SoC stages shared by both flows: acceleration-target handling, accelerator implementation, functional verification, HLS synthesis, and ESP integration. FPGA board execution of the SoCPilot designs is reported separately in RQ1 using the manually prepared runtime code and is not used to give SoCPilot an advantage in the agent comparison.

\textbf{Area-constrained design evaluation.}
For RQ3, we vary the application-specific area constraint and examine how SoCPilot adapts the generated architecture. For each constraint, we record the selected CPU, accelerator set, relevant system parameters, estimated area, and resulting end-to-end latency. The FPGA board is used to validate the selected configurations, whereas area is obtained from the common analytical area model rather than from FPGA resource utilization.

A generated SoCPilot design is counted as an \emph{end-to-end validated design} when it produces a valid SoC configuration, synthesizable accelerators, a successfully generated FPGA bitstream, correct accelerator invocation through the manually prepared embedded runtime, and application outputs that pass the predefined correctness checks against the CPU reference.

\subsubsection{Application Benchmark}
\label{subsec:application_suite}
Table~\ref{tab:e2e_applications} summarizes five end-to-end edge benchmarks used in our evaluation: license-plate processing (\texttt{plate\_recognition}), color-based image recognition (\texttt{recognition}), aerial video analytics (\texttt{wami}), document image processing (\texttt{scan}), and low-light vision enhancement (\texttt{nightvision}). Each benchmark is implemented as a complete C/C++ application rather than an isolated kernel, preserving application-level control flow and multiple dependent processing stages. The candidate functions exhibit different computational intensity, memory-access behavior, and acceleration potential, making the suite suitable for evaluating end-to-end CPU--accelerator SoC design across diverse application characteristics.

\begin{table*}[t]
\centering
\caption{End-to-end edge-application suite. The hotspot column is filled with the functions selected by SoCPilot after profiling and cost evaluation.}
\label{tab:e2e_applications}
\renewcommand{\arraystretch}{1.12}
\resizebox{\textwidth}{!}{
\begin{tabular}{p{2.7cm}p{2.55cm}p{4.85cm}p{4.05cm}p{3.15cm}}
\toprule
\textbf{Application} &
\textbf{Scenario} &
\textbf{Major processing stages} &
\textbf{Application-level quality metric} &
\textbf{Selected hotspot(s)} \\
\midrule
License-plate recognition (\texttt{plate\_recognition}) &
Roadside or parking camera &
Image loading, grayscale conversion, morphological enhancement, edge/region extraction, candidate filtering, and plate localization &
Detection success / bounding-box consistency &
\texttt{morphology\_tophat}, \texttt{morphology\_blackhat} \\

Color-based object recognition (\texttt{recognition}) &
Embedded visual recognition &
BGR-to-HSV conversion, color mask generation, mask dilation, contour localization, and target-center estimation &
Recognition accuracy / localization error &
\texttt{bgr\_to\_hsv}, \texttt{dilate\_mask} \\

Document scanning (\texttt{scan}) &
Mobile or edge document scanner &
RGB-to-gray conversion, adaptive thresholding, morphology, edge extraction, document mask generation, and feature merging &
Document mask accuracy / binarization quality &
\texttt{scan\_adaptive\_threshold} \\

Wide-area motion imagery (WAMI) &
Aerial video analytics &
Image debayering, geometric warping, gradient-based Lucas--Kanade tracking, Hessian construction, and foreground modeling &
Tracking error / foreground detection quality &
\texttt{warp\_image}, \texttt{steepest\_descent}, \texttt{hessian} \\

Night-vision enhancement (\texttt{nightvision}) &
Low-light smart camera &
Filtering, median denoising, histogram equalization, convolutional enhancement, wavelet transform, and image fusion &
Enhanced-image quality / contrast preservation &
\texttt{conv2d} \\
\bottomrule
\end{tabular}}
\end{table*}

